% !TeX root = ../../Book2.tex
% 纲要标题：### 12.1 映射和截面的延拓问题
% 纲要位置：几何学计划纲要.md，第 1472 行。
% 正文状态：初稿；尚未正式审核。

% 纲要内容开始
% 
% * 空间对、骨架与相对条件
% * 单个胞腔上的延拓判据
% * 基本群作用产生的系数系统
% 
% 纲要内容结束

\section{映射和截面的延拓问题}
\label{g2:sec:1201}

给圆周上的每一点指定一个单位向量以后，能否把这个选择连续地填到整个圆盘？如果向量已经绕了一整圈，填充时便会遇到障碍。更高维的情形保留了同一结构：一个新胞腔是圆盘，旧数据已经规定它的边界；是否能够填进去，由边界在目标或纤维中的同伦类决定。先把这些逐胞腔的问题写准确，下一节才能把它们组织为上同调类。

\subsection{空间对、骨架与相对条件}
\label{g2:subsec:1201:01}

设 \((X,A)\) 为 CW 对，给定映射 \(f_A:A\to Y\)。\term[yingshe-yantuo-wenti]{映射延拓问题}[extension problem]要求寻找 \(f:X\to Y\)，满足 \(f|_A=f_A\)。这里边界数据是实际给定的映射；若只要求其同伦类相同，就需要另作说明。记
\begin{equation}\label{g2:eq:relative-obstruction-skeleta}
	X_{-1}=A,\qquad X_r=A\cup X^r\quad(r\ge0).
\end{equation}
这些是相对骨架：即使 \(A\) 含高维胞腔，它也从一开始就全部保留。下标 \(r\) 只控制 \(X\setminus A\) 中已经加入的胞腔；它与上一章表示同伦截段的 \(X[n]\) 含义不同。

\ProseParagraphBreak

更一般地，给定纤维化 \(p:E\to X\) 及 \(A\) 上的截面 \(s_A\)，要求寻找满足
\[
	p\circ s=\operatorname{id}_X,\qquad s|_A=s_A
\]
的截面。映射延拓是它的特例：取 \(E=X\times Y\)，截面必形如 \(s(x)=(x,f(x))\)。而对一般提升问题，若已有 \(g:X\to B\)、纤维化 \(p_0:E_0\to B\) 和 \(g|_A\) 的一个指定提升，则先取拉回
\[
	g^*E_0=\{(x,e):g(x)=p_0(e)\}\longrightarrow X.
\]
截面 \(x\mapsto(x,\widetilde g(x))\) 与满足原边界条件的提升 \(\widetilde g\) 一一对应。于是映射、提升和截面可以用同一套逐骨架方法处理。

\ProseParagraphBreak

从 \(X_r\) 到 \(X_{r+1}\) 时，每个新胞腔的边界已经有值，同维胞腔内部则互不相交。只要对每个特征圆盘找到保持边界的延拓，就能沿原来的附着映射粘合。CW 弱拓扑保证最后所得映射连续：连续性只需在 \(A\) 及每个特征圆盘上核对。对无限 CW 复形，这个事实允许逐维构造相容映射；在各个有限阶段分别找到互不相容的延拓，则还没有得到一张定义在 \(X\) 上的映射。

\begin{example}[圆盘边界上的绕数]\label{g2:exm:boundary-winding-extension}
	令 \(X=D^2\)、\(A=S^1\)，取 \(f_A(z)=z^k:S^1\to S^1\)，其中 \(k\in\mathbb Z\)。若它延拓到 \(D^2\)，把圆盘沿半径收缩到中心，就得到 \(f_A\) 的自由零伦，故绕数必须为零。反过来，\(k=0\) 时常值映射本身便给出延拓。因此延拓恰在 \(k=0\) 时存在。这同时是平凡实二维丛上的单位截面延拓问题：单位向量的边界绕数阻止它无零点地填入圆盘。
\end{example}

还须区分两种要求：把一个已经选好的 \(s_r:X_r\to E\) 原样延到下一层，和只固定 \(X_{r-1}\) 上的数据、允许重新选择 \(r\) 胞腔内的截面。第一种要求逐胞腔的障碍值全部为零；第二种允许用新选择改变这些值，因而将在下一节产生“障碍类为零”的判据。

\ProseParagraphBreak

比较两张延拓的同伦也属于同一问题。若 \(s_0,s_1\) 在 \(A\) 上相同，寻找固定 \(A\) 的截面同伦，就是在
\[
	\bigl(X\times I,\ (X\times\partial I)\cup(A\times I)\bigr)
\]
上延拓截面，其中两端取 \(s_0,s_1\)，侧面保持给定值。每个 \(r\) 胞腔乘区间成为一个 \((r+1)\) 维延拓问题，这将解释差异类的次数为何比存在性障碍低一度。

\subsection{单个胞腔上的延拓判据}
\label{g2:subsec:1201:02}

先看普通映射。设 \(q\ge1\)，\(\chi:D^{q+1}\to X_{q+1}\) 是一个新胞腔的特征映射，\(f_q:X_q\to Y\) 已经给定。复合
\[
	S^q\xrightarrow{\chi|_{S^q}}X_q\xrightarrow{f_q}Y
\]
就是待填圆盘上的边界数据。

\begin{proposition}[单个胞腔的映射延拓]\label{g2:prop:cell-extension-criterion}
	设 \(q\ge1\)，\(Y\) 为道路连通空间。映射 \(u:S^q\to Y\) 能延拓到 \(D^{q+1}\)，当且仅当它自由零伦；选择球面基点及目标基点识别后，这等价于对应的 \(\pi_q(Y)\) 类为零。因此，\(f_q:X_q\to Y\) 能保持已有值延到 \(X_{q+1}\)，当且仅当每个相对 \((q+1)\) 胞腔给出的这个类均为零。
\end{proposition}
\begin{proof}
	若 \(U:D^{q+1}\to Y\) 延拓 \(u\)，则 \(H(z,t)=U((1-t)z)\) 把边界映射收缩到 \(U(0)\)。反过来，给定从 \(u\) 到常值的同伦，将 \(D^{q+1}\) 写为 \(S^q\) 的锥，这个同伦就在锥顶处下降为所需延拓。

	由\cref{g2:cor:free-homotopy-orbits}，自由同伦对应基本群对带基点同伦群的作用轨道。任何作用自同构都保持零元，且只有零元能被送到零元，因此自由零伦等价于任一基点识别下的类为零。最后逐胞腔作上述填充；它们在公共边界均等于 \(f_q\)，所以可粘合为下一相对骨架上的连续映射。
\end{proof}

零胞腔要求目标非空；一胞腔要求它的两个端点像可以用道路相连。这两步是非空性和道路分支的问题。二胞腔开始才出现 \(\pi_1\) 中的边界环路；若这个群非交换，仍可判定单个环路是否为单位元，但尚不能把全部障碍当作普通阿贝尔群系数的上链。

\ProseParagraphBreak

截面的圆盘判据需要把边界数据送到同一纤维。设 \(p:E\to X\) 为 Hurewicz 纤维化，拉回到特征圆盘得到 \(\pi:V=\chi^*E\to D^{q+1}\)。第11章的\cref{g2:lem:fibration-disk-model}给出保持基底的纤维同伦等价
\[
	V\ \underset{P}{\stackrel{Q}{\rightleftarrows}}\ D^{q+1}\times F,
\]
其中 \(F=\pi^{-1}(0)\)。给定的边界截面 \(s_\partial:S^q\to V\) 经 \(Q\) 写成 \(Q s_\partial(z)=(z,u(z))\)。由此得到边界的纤维映射 \(u:S^q\to F\)。

\begin{proposition}[单个胞腔的截面延拓]\label{g2:prop:section-cell-extension}
	设 \(q\ge1\)，\(\pi:V\to D^{q+1}\) 为具有非空道路连通纤维的 Hurewicz 纤维化，\(s_\partial\) 为其边界上的截面。在上述纤维同伦乘积模型中，\(s_\partial\) 可延拓为整个圆盘上的截面，当且仅当 \(u:S^q\to F\) 的类为零。这个零值判据与乘积模型及基点识别的选择无关。
\end{proposition}
\begin{proof}
	若截面 \(s\) 已存在，\(Qs\) 的纤维坐标便延拓 \(u\)，所以其类为零。

	反过来，假定 \(u\) 延拓为 \(U:D^{q+1}\to F\)。复合 \(t(x)=P(x,U(x))\) 是圆盘上的截面；其边界未必等于 \(s_\partial\)，但 \(PQ\) 与恒等映射之间的纤维同伦给出从 \(t|_{S^q}\) 到 \(s_\partial\) 的截面同伦。对 CW 对 \((D^{q+1},S^q)\) 应用\cref{g2:lem:relative-homotopy-lift}，基底同伦取 \(H(x,r)=x\)，初始提升取 \(t\)，边界提升取刚才的同伦。所得提升的终时刻是一个严格延拓 \(s_\partial\) 的截面。这一步把纤维同伦模型给出的边界相等到同伦，修正成实际边界相等。

	判据两端已经等价于与模型无关的截面存在性，因此改变模型不改变它是否为零；基点改变也保持零元。
\end{proof}

纤维丛在圆盘上可用实际平凡化完成同样的论证；保留纤维同伦模型则同时适用于路径纤维化。对一般 Serre 纤维化，相同的圆盘论证可以只提升有限 CW 参数：沿径向收缩提升已给边界，得到中心纤维中的球面映射；若它零伦，就把这张圆盘填充作为初始提升，再反向提升径向同伦，并在边界指定原提升的反向过程。相对提升引理保证最终边界恰为原截面。这给出同一零值判据，而不要求整条纤维作为任意参数空间参与提升。

\begin{example}[连通度给出的自动延拓]\label{g2:exm:connectivity-extension-range}
	设 \(n\ge2\)，\(Y\) 为 \((n-1)\) 连通空间，\((X,A)\) 为 CW 对，\(f_A:A\to Y\) 已给定。选择新顶点的像，并用道路填入新一胞腔；随后 \(\pi_q(Y)=0\)（\(1\le q<n\)）使所有维数不超过 \(n\) 的新胞腔都可填入。因此 \(f_A\) 总可延到 \(X_n\)，首次可能出现的障碍位于 \((n+1)\) 胞腔，取值于 \(\pi_n(Y)\)。对具有 \((n-1)\) 连通非空纤维的纤维化，指定截面的相对延拓具有相同范围。
\end{example}

\subsection{基本群作用产生的系数系统}
\label{g2:subsec:1201:03}

单个胞腔只问一个球面类是否为零，所以任意基点运输都不会改变答案。若要把多个胞腔上的障碍相加，或比较两次选择造成的差，便必须说明这些类处于哪些群、不同群之间怎样识别。第5章的局部系统在这里重新出现，系数从纤维同调换成纤维同伦群。

\begin{definition}[简单空间]\label{g2:def:simple-space}
	道路连通空间 \(Y\) 若满足 \(\pi_1(Y)\) 交换，且它对每个 \(\pi_q(Y)\)（\(q\ge2\)）的基点运输作用都平凡，称为\term[jiandan-kongjian]{简单空间}[simple space]。对整数 \(q\ge1\)，若只要求这些作用在 \(2\le j\le q\) 时平凡，则称为 \(q\) 简单空间。\(q=1\) 时这个条件仅要求基本群交换。
\end{definition}

基本群对自身的作用是共轭，因此定义中的交换性恰好使一次群的基点运输也不受路径选择影响。由\cref{g2:prop:basepoint-action}，两条具有相同端点的道路给出的基点变换相差一个环路作用；对简单空间，这个差别消失。于是其各点的 \(\pi_q\) 可以作唯一的路径无关识别。每个单连通空间都是简单空间；\(S^1\) 也简单，因为其基本群为 \(\mathbb Z\)，高阶同伦群全为零。

\begin{proposition}[纤维同伦群的局部系统]\label{g2:prop:homotopy-local-system}
	设 \(p:E\to X\) 为基底道路连通、纤维非空且简单的 Hurewicz 纤维化。对每个 \(q\ge1\)，纤维同伦群组成阿贝尔群的局部系统 \(\Pi_q(p)\)：在 \(x\in X\) 处取群 \(\pi_q(F_x)\)，沿道路 \(\gamma:x\to y\) 的同态由纤维运输 \(T_\gamma:F_x\to F_y\) 诱导，并用纤维内部的路径识别基点。它只依赖于 \(\gamma\) 的定端点同伦类，且
	\begin{equation}\label{g2:eq:homotopy-local-transport}
		\Pi_q(p)(\gamma*\eta)=\Pi_q(p)(\eta)\Pi_q(p)(\gamma).
	\end{equation}
	若只处理次数至 \(q\)，纤维 \(q\) 简单已足够。
\end{proposition}
\begin{proof}
	第11章的道路提升把 \(T_\gamma\) 确定到自由同伦；反向道路给出同伦逆，定端点道路同伦也给出相应映射同伦。若为 \(F_x,F_y\) 各选一点，\(T_\gamma\) 的基点像未必等于已选终点基点，故还须沿 \(F_y\) 内的一条道路变换基点。简单性说明这样的基点变换与所选道路无关。改变 \(T_\gamma\) 的提升选择时，同伦的基点轨迹也只造成同一类已经消失的环路作用，因而群同态无歧义。

	纤维运输的拼接满足 \(T_{\gamma*\eta}\simeq T_\eta T_\gamma\)，诱导同态遂满足\eqref{g2:eq:homotopy-local-transport}；常值道路诱导恒等映射。这正是\cref{g2:def:local-system}的局部系统公理。\(q\ge2\) 时同伦群本来就交换，\(q=1\) 时交换性由简单性提供。对 Serre 纤维化，可通过第11章的 Hurewicz 替换及纤维弱同伦等价搬回这一系统，或直接对球面参数作上述运输。
\end{proof}

“纤维简单”消除的是同一纤维内的基点歧义，并不要求绕基底一周后的运输恒等。例如取反足映射 \(a:S^2\to S^2\)，其映射环面
\[
	E=(S^2\times I)/((z,1)\sim(a(z),0))\longrightarrow S^1
\]
具有单连通纤维，但基底正向绕行在 \(\pi_2(S^2)=\mathbb Z\) 上作用为乘 \(-1\)。这是因为 Hurewicz 同构把该作用对应到反足映射在二次同调上的度 \(-1\)。因此 \(\Pi_2(p)\) 是反号局部系统，不能以常值整数系数取代。定向 \(S^m\) 丛的情况则相反：其顶次运输保持定向，故 \(\Pi_m(p)\cong\mathbb Z\) 为常值系统；未定向时，它记录的正是纤维定向反转。

\ProseParagraphBreak

普通映射问题中，若 \(Y\) 简单，乘积纤维化 \(X\times Y\to X\) 的运输为恒等，因而系数就是常值群 \(\pi_q(Y)\)。若目标不简单，则已给的低维映射也会进入系数。例如一张定义在 \(X_2\) 上的映射决定
\[
	\rho:\pi_1(X)\longrightarrow\pi_1(Y),
\]
因为 \(X\) 的基本群由二骨架决定。固定 \(\rho\) 后，\(\pi_1(Y)\) 在 \(\pi_q(Y)\) 上的作用可沿 \(\rho\) 拉回，成为 \(X\) 上的局部系统。仅给出抽象群 \(\pi_q(Y)\)，还没有指定这组系数。下一节采用简单纤维而允许基底运输不平凡的条件，既保持阿贝尔上同调的运算，也包括非定向球丛这样的几何情形。

\ProseParagraphBreak

因此，第 \(q+1\) 层的障碍将由相对胞腔上链记录在
\[
	C^{q+1}(X,A;\Pi_q(p))
\]
中。每个胞腔的定向决定球面类的符号，道路运输决定不同胞腔的系数如何相加。下一节要证明这些值满足余圈条件，并确定修改已有延拓时产生的余边界；只有完成这两步，才能把逐胞腔判据提升为上同调中的不变量。映射与提升问题的阅读对照可参看\cite[第4.3节，Obstruction Theory]{Hatcher2002}及\cite[第18章，第5节，Obstruction theory]{May1999}。
% 纲要条目排为编号小节标题，并进入目录与书签；说明段仍保留于纲要注释中。
